\begin{table}[H]
\small
\setlength{\tabcolsep}{4.2pt}
\renewcommand{\arraystretch}{0.95}

\caption{\label{tab:ea20-income-gap-product-contributions}Product-group contributions to euro-area annual income inflation gaps, 2021--2024}
\centering
\begin{tabular}[t]{llll}
\toprule
Product group & Q1 contribution & Q5 contribution & Q1-Q5 contribution gap\\
\midrule
Food and non-alcoholic beverages & 1.3 & 1.1 & 0.2\\
Alcoholic beverage, tobacco and narcotics & 0.4 & 0.2 & 0.2\\
Clothing and footwear & 0.1 & 0.2 & -0.0\\
Housing (rentals and repairs) & 0.2 & 0.2 & 0.0\\
Water, electricity, gas and other fuels & 1.3 & 1.0 & 0.3\\
\addlinespace
Transport & 0.4 & 0.7 & -0.3\\
Other & 0.1 & 0.2 & -0.1\\
\midrule
\textbf{Total} & \textbf{4.0} & \textbf{3.6} & \textbf{0.4}\\
\bottomrule
\end{tabular}
\vspace{-0.4em}
\begin{flushleft}
\footnotesize{Notes. Entries are percentage points. Contributions are annual-average monthly inflation contributions summed over 2021--2024. Product groups use the default mapping in \texttt{plot\_contribution\_gap}. France is computed at COICOP level 3 with INSEE income quintiles; the euro area is rebuilt from country-level contributions and EA20 HICP country weights. The gap is Q1 minus Q5. This table decomposes annual contribution gaps, not the cumulative price-index gap in Table~\ref{tab:country-heterogeneity-cumulative-gap}. Sources: Eurostat, authors' calculations.}
\end{flushleft}
\end{table}
